%----------------------------------------------------------------------------
\chapter{Methodology and implementation}
\label{chap:Methodology}
%----------------------------------------------------------------------------

This chapter describes the implementation that produces the data in
\autoref{chap:Experiments} and, more importantly, the decisions that keep the
XPBD-vs-VBD comparison fair.

%----------------------------------------------------------------------------
\section{Platform and environment}
\label{sec:Platform}
%----------------------------------------------------------------------------
Both simulators are implemented inside the Unity Editor, and all measurements
reported in \autoref{chap:Experiments} were collected in the same environment:
\begin{itemize}
    \item \textbf{CPU:} AMD Ryzen 7 5800H
    \item \textbf{Memory:} 16~GB DDR4
    \item \textbf{GPU:} NVIDIA GeForce RTX 3060 (Laptop)
    \item \textbf{Operating system:} Windows 11 Pro
    \item \textbf{Development tool:} Unity 6000.4.0f1
\end{itemize}

Several of Unity's built-in tools proved useful and easy to work with while
conducting the experiments: the Profiler for gathering FPS statistics, the
Recorder for capturing videos of the simulations, and Shader Graph for creating
textures in a straightforward way.

%----------------------------------------------------------------------------
\section{Integration architecture}
\label{sec:Architecture}
%----------------------------------------------------------------------------
Both solvers share a common scene scaffolding, and the simulations are driven by
the \texttt{Update} method of Unity's \texttt{MonoBehaviour} classes. The
scaffolding handles mesh and spring-topology construction at scene load, the
recording of statistics such as frame time and memory, and rendering, which runs
on a separate thread from the simulator.

In the Unity Editor, the \texttt{Update} method is called as often as the
environment allows; if a step takes more compute time than the frame time, the
next frame is delayed until it finishes. The \texttt{Update} method therefore
effectively provides a variable frame rate that can be used to measure
computation time. To keep the experiments deterministic and easier to control, I
use a fixed simulation time step of $dt = 1/30$~s. This is also a relatively large
time step, which makes it a good test of the solvers' stability and lets the
substep count carry more meaning in \autoref{ssec:ErrorVsMs}.

\paragraph{Overhead of the shared scaffolding.}
The scaffolding adds a small, practically constant per-frame overhead, mostly from
Unity-specific calls and rendering. I measure FPS with the Unity Profiler, which
counts the actual frames drawn, so the reported value includes this overhead; but
because both solvers pay it identically, it does not bias the comparison. To keep
performance as high as possible, I favor primitive types, arrays, and structs over
classes and lists, and I cache frequently accessed variables globally to shift the
load from the CPU to memory.

\subsection{Cloth model construction}
\label{ssec:ClothModel}
The mass-spring cloth model I implement is the one used by
M\"uller~\cite{tenMinutePhysics}. Each internal vertex is connected to its
neighbors by three types of distance constraint, each with its own
stiffness/compliance value. \textbf{Stretching} constraints connect a vertex to
its immediate horizontal and vertical neighbors; \textbf{shear} constraints
connect it to its diagonal neighbors; and \textbf{bending} constraints connect it
to its second horizontal and vertical neighbors. An internal vertex therefore has
twelve constraints attached to it. Note that the bending constraints are not
dihedral-angle constraints: I deliberately use a single constraint form to keep
the comparison less complex.

The mass of each vertex is set proportionally to the area of cloth it represents:
a vertex receives a quarter of the summed area of the quadrilaterals it touches,
and the parameter the user supplies is the total mass of the cloth. This keeps the
cloth's behavior consistent across resolutions. Assigning the same mass to every
vertex instead would also leave the side and corner vertices effectively
heavier relative to the area they represent.

%----------------------------------------------------------------------------
\section{How the comparison is kept fair}
\label{sec:Fairness}
%----------------------------------------------------------------------------
I adopt several explicit fairness controls, described below.

\subsection{Stiffness $\leftrightarrow$ compliance mapping}
\label{ssec:StiffCompl}
VBD takes a stiffness $k$ directly, whereas XPBD takes a compliance
$\alpha = 1/k$. To match materials, I set
\begin{equation}
    \alpha_{\mathrm{XPBD}} \;=\; 1/k_{\mathrm{VBD}}.
\end{equation}
This mapping is verified in \autoref{ssec:HarmonicOscillator}: a single spring
with known analytic parameters yields the same oscillation frequency and decay
rate in both solvers.

\subsection{Matched compute budget --- and the substep/iteration caveat}
\label{ssec:ComputeBudget}
The natural unit of work differs between the two methods: XPBD's primary lever is
the substep count, while VBD's is the inner iteration count. \textbf{A substep is
not an iteration.} A substep recomputes the inertial target $\mathbf{y}$ and
solves a fresh constraint system, whereas an iteration refines the same step.
Because of this asymmetry, I compare the two methods on the basis of
\emph{wall-clock time} throughout \autoref{chap:Experiments}. One source of bias
is that VBD uses a maximum iteration count while XPBD uses a fixed iteration
count. Either method could implement either scheme, but these choices are
inherent to each method's design (even allowing for their future optimizations),
so I keep them as they are.

\subsection{Matched gravity, $dt$, and damping}
\label{ssec:MatchedDtGravity}
Unless otherwise noted, all experiments use gravity
$\mathbf{g} = (0, -9.81, 0)$~m/s$^{2}$ and a simulation time step of
$dt = 1/30$~s. Rayleigh damping is disabled in both methods, because the
difference in solver structure produces different results from the same damping
parameters.

\subsection{Handling VBD's accelerations for the core comparison}
\label{ssec:DisableChebyshev}
The 2024 VBD paper~\cite{Chen2024} reports significant speedups from Chebyshev
acceleration and adaptive initialization. Adaptive initialization is VBD's own
inertial prediction, which is also present in XPBD (which guesses with explicit
Euler), so I keep it enabled. Chebyshev acceleration is likewise part of VBD, but
it confounds an apples-to-apples comparison, since XPBD could implement the same
technique, as shown in \cite{10.1145/2816795.2818063}; I therefore disable it for
the core comparison.

\subsection{Hessian projection for stability}
\label{ssec:HessianProjection}
The original VBD paper argues that the Hessian need not be positive-definite: even
if a step occasionally does not point in a descent direction, the small number of
such steps does not affect overall convergence. In my experiments, however,
projecting the matrix to be positive-definite let the method handle more extreme
scenarios without noticeably changing its behavior or performance, so I project
the Hessian to be positive-definite.

%----------------------------------------------------------------------------
\section{Self-collision handling}
\label{sec:SelfCollision}
%----------------------------------------------------------------------------
The two solvers implement self-collision differently, by necessity.
\emph{XPBD self-collision} uses non-penetration constraints projected each
iteration --- contact is enforced as a position constraint. \emph{VBD
self-collision}, in contrast, uses an elastic penalty potential added to $E$,
contributing a gradient and Hessian to the affected vertex blocks.

This is a model difference, not a solver difference: a penalty model \emph{could}
be implemented in XPBD and a projection model \emph{could} be implemented in VBD,
but the natural, most solver-fitting form differs between them. The self-collision
experiment in \autoref{ssec:SelfCollisionDrop} is therefore purely qualitative.

%----------------------------------------------------------------------------
\section{Parallelization (forward-looking)}
\label{sec:Parallelisation}
\label{ssec:Parallelisation}
%----------------------------------------------------------------------------
All measurements in this thesis are run single-threaded on the CPU. As noted here
and revisited in \autoref{sec:ThreatsValidity}, the two solvers have different
parallelization profiles. \emph{VBD} parallelizes naturally over vertex blocks
through graph coloring: vertices that share no constraint can be updated
independently within a color class, which suits GPU execution well. \emph{XPBD}
can be parallelized through edge coloring, but this coloring is finer-grained
(per constraint --- of which there are usually many more than vertices --- rather
than per vertex), so it produces more color classes and hence more sequential
dependencies.

The expected result --- that parallelization favors VBD more than XPBD --- is
plausible but unverified by my experiments. It stands as a threat to the external
validity of any claim about which solver is ``faster overall.''
